\documentclass[11pt,a4paper]{article}

\usepackage{fontspec}
\usepackage{polyglossia}
\setmainlanguage{russian}
\setmainfont{Liberation Serif}
\setsansfont{Liberation Sans}
\setmonofont{DejaVu Sans Mono}

\usepackage{amsmath,amssymb}
\usepackage{booktabs}
\usepackage{array}
\usepackage{geometry}
\geometry{margin=1.8cm}

\setlength{\parskip}{0.45em}
\setlength{\parindent}{0pt}

\newcommand{\num}[1]{\(#1\)}

\begin{document}

\begin{center}
{\Large Ручной расчёт одной строки}\\[0.4em]
{\large Лабораторная работа №4, вариант 2}\\[0.2em]
Метод средних прямоугольников с удвоением числа шагов
\end{center}

Сверяем программу \texttt{3.py} с прямым счётом для одной точки таблицы.
Берём строку \(j = 7\). Это первая строка, где четырёх шагов уже не хватает
и программа переходит к восьми. На ней видно и саму формулу, и правило остановки.

\section*{Что считаем}

Параметры варианта~2:
\[
a = 0,\quad b = 2,\quad c = 0,\quad d = 1,\quad m = 20,\quad \varepsilon = 0{,}001,
\]
\[
f(x, t) = \cos\left(\frac{t}{1 + x^{2}} + 0{,}001\, x\right).
\]

Шаг по параметру \(t\):
\[
\tau = \frac{d - c}{m} = \frac{1}{20} = 0{,}05,
\qquad
t_{7} = c + 7\tau = 0{,}35.
\]

Нужна площадь
\[
F(0{,}35) = \int_{0}^{2} f(x,\, 0{,}35)\, dx.
\]

Отрезок \([0, 2]\) делим на \(N\) равных частей длины \(h = 2/N\).
На \(i\)-й части (нумерация с нуля) берём середину
\(x_{i} = (i + 0{,}5)\, h\) и заменяем график прямоугольником этой высоты:
\[
S_{N} = h \sum_{i=0}^{N-1} f(x_{i},\, 0{,}35).
\]

Считаем \(S_{2}\), затем \(S_{4}\), затем \(S_{8}\).
Останавливаемся, когда соседние оценки впервые отличаются меньше чем на \(0{,}001\).
В ответ берём более мелкую из этих двух оценок. Так же делает программа.

Значения косинуса округлены до 10 знаков после запятой.
Аргумент косинуса везде в радианах.

\section*{Шаг 1. Два прямоугольника, \(N = 2\)}

\[
h = \frac{2}{2} = 1.
\]

Середины кусков \([0, 1]\) и \([1, 2]\): \(x = 0{,}5\) и \(x = 1{,}5\).

\paragraph{Точка \(x = 0{,}5\).}
\[
1 + x^{2} = 1 + 0{,}25 = 1{,}25,
\qquad
\frac{0{,}35}{1{,}25} = 0{,}28,
\qquad
0{,}001 \cdot 0{,}5 = 0{,}0005.
\]
Аргумент: \(0{,}28 + 0{,}0005 = 0{,}2805\).
\[
f(0{,}5) = \cos 0{,}2805 = 0{,}9609171404.
\]

\paragraph{Точка \(x = 1{,}5\).}
\[
1 + x^{2} = 1 + 2{,}25 = 3{,}25,
\qquad
\frac{0{,}35}{3{,}25} = 0{,}1076923077,
\qquad
0{,}001 \cdot 1{,}5 = 0{,}0015.
\]
Аргумент: \(0{,}1076923077 + 0{,}0015 = 0{,}1091923077\).
\[
f(1{,}5) = \cos 0{,}1091923077 = 0{,}9940444408.
\]

Сумма высот и площадь:
\[
0{,}9609171404 + 0{,}9940444408 = 1{,}9549615812,
\]
\[
S_{2} = 1 \cdot 1{,}9549615812 = 1{,}9549615812.
\]

Сравнивать пока не с чем: это первая сетка. Запоминаем \(S_{2}\) и удваиваем число шагов.

\section*{Шаг 2. Четыре прямоугольника, \(N = 4\)}

\[
h = \frac{2}{4} = 0{,}5.
\]

Середины: \(0{,}25;\ 0{,}75;\ 1{,}25;\ 1{,}75\).

{\small
\begin{center}
\begin{tabular}{@{}r r r r r r@{}}
\toprule
\(x\) & \(1+x^{2}\) & \(0{,}35/(1+x^{2})\) & \(0{,}001\,x\) & аргумент & \(f(x)\) \\
\midrule
0,25 & 1,0625 & 0,3294117647 & 0,00025 & 0,3296617647 & 0,9461518922 \\
0,75 & 1,5625 & 0,224 & 0,00075 & 0,22475 & 0,9748498532 \\
1,25 & 2,5625 & 0,1365853659 & 0,00125 & 0,1378353659 & 0,9905157359 \\
1,75 & 4,0625 & 0,0861538462 & 0,00175 & 0,0879038462 & 0,9961389441 \\
\bottomrule
\end{tabular}
\end{center}
}

Проверка одной строки, \(x = 0{,}75\):
\[
x^{2} = 0{,}5625,\quad 1 + x^{2} = 1{,}5625,\quad \frac{0{,}35}{1{,}5625} = 0{,}224,
\]
\[
0{,}224 + 0{,}00075 = 0{,}22475,\quad \cos 0{,}22475 = 0{,}9748498532.
\]

Сумма четырёх значений:
\[
0{,}9461518922 + 0{,}9748498532 + 0{,}9905157359 + 0{,}9961389441 = 3{,}9076564254.
\]
Площадь:
\[
S_{4} = 0{,}5 \cdot 3{,}9076564254 = 1{,}9538282127.
\]

Сравнение с предыдущей сеткой:
\[
|S_{4} - S_{2}| = |1{,}9538282127 - 1{,}9549615812| = 0{,}0011333685.
\]
\(0{,}0011333685 > 0{,}001\). Заданная точность ещё не достигнута, число шагов снова удваиваем.

\section*{Шаг 3. Восемь прямоугольников, \(N = 8\)}

\[
h = \frac{2}{8} = 0{,}25.
\]

Середины: \(0{,}125;\ 0{,}375;\ 0{,}625;\ 0{,}875;\ 1{,}125;\ 1{,}375;\ 1{,}625;\ 1{,}875\).

{\small
\begin{center}
\begin{tabular}{@{}r r r@{}}
\toprule
\(x\) & аргумент & \(f(x)\) \\
\midrule
0,125 & 0,3447403846 & 0,9411632220 \\
0,375 & 0,3072243151 & 0,9531766469 \\
0,625 & 0,2523103933 & 0,9683382358 \\
0,875 & 0,1991050885 & 0,9802439768 \\
1,125 & 0,1556077586 & 0,9879175225 \\
1,375 & 0,1224560811 & 0,9925116188 \\
1,625 & 0,0977623391 & 0,9952250674 \\
1,875 & 0,0793836505 & 0,9968507723 \\
\bottomrule
\end{tabular}
\end{center}
}

Проверка одной строки, \(x = 0{,}125\):
\[
x^{2} = 0{,}015625,\quad 1 + x^{2} = 1{,}015625,
\]
\[
\frac{0{,}35}{1{,}015625} = 0{,}3446153846,
\qquad
0{,}001 \cdot 0{,}125 = 0{,}000125,
\]
\[
0{,}3446153846 + 0{,}000125 = 0{,}3447403846,
\]
\[
\cos 0{,}3447403846 = 0{,}9411632220.
\]

Сумма восьми значений:
\begin{align*}
&0{,}9411632220 + 0{,}9531766469 + 0{,}9683382358 + 0{,}9802439768 \\
&+ 0{,}9879175225 + 0{,}9925116188 + 0{,}9952250674 + 0{,}9968507723 \\
&= 7{,}8154270624.
\end{align*}

Площадь:
\[
S_{8} = 0{,}25 \cdot 7{,}8154270624 = 1{,}9538567656.
\]

Сравнение с сеткой из четырёх шагов:
\[
|S_{8} - S_{4}| = |1{,}9538567656 - 1{,}9538282127| = 0{,}0000285529.
\]
\(0{,}0000285529 < 0{,}001\). Критерий выполнен, дальше удваивать не нужно.

\section*{Ответ и сверка с программой}

\[
F(0{,}35) \approx S_{8} = 1{,}95385677, \qquad N = 8.
\]

В таблицу пишется \(S_{8}\), а не \(S_{4}\): берётся сумма на той сетке, где условие уже выполнилось.
Восемь знаков после запятой~--- это округление \(1{,}9538567656\). В таком виде число печатает программа:

\begin{center}
\texttt{7 | 0.3500 | 1.95385677 | 8}
\end{center}

Ручной счёт даёт те же \(t_{7} = 0{,}35\), \(F = 1{,}95385677\) и \(N = 8\).

Почему именно 8, а не 4: разность \(S_{4}\) и \(S_{2}\) ещё больше допуска \(0{,}001\)
(получилось \(0{,}00113\)). Разность \(S_{8}\) и \(S_{4}\) уже меньше допуска
(получилось \(0{,}000029\)).

\end{document}
